\subsection{Bối cảnh và mục tiêu}
Trong bãi xe ngoài trời, việc xác định vị trí còn trống từ camera có thể hỗ trợ người quản lý theo dõi mức sử dụng bãi mà không cần kiểm tra từng ô trực tiếp. Bài toán gặp khó khăn khi xe ở xa camera, bóng đổ che đường kẻ, ánh sáng thay đổi hoặc mặt đường bị ảnh hưởng bởi mưa.

Đề tài xây dựng hệ thống nhận diện trạng thái ô đỗ và đếm số chỗ trống từ ảnh màu toàn cảnh. Hệ thống được thiết kế để so sánh ba phương pháp: phân loại ảnh từng ô bằng HOG kết hợp SVM, phát hiện xe bằng YOLO11n trên toàn ảnh và phát hiện xe bằng cùng detector trên các ROI/tile. Việc so sánh xem xét chất lượng nhận diện, sai số đếm và chi phí xử lý; lợi ích của chia tile cần được kiểm chứng bằng số liệu.

\subsection{Đầu vào và đầu ra}
Đầu vào gồm ảnh màu toàn cảnh từ camera cố định và cấu hình các ô đỗ tương ứng với ảnh. Mỗi ô có mã định danh và polygon biểu diễn vùng ô trong hệ tọa độ ảnh gốc. Ngày thu nhận, camera và thời tiết được giữ lại để chia dữ liệu và phân tích kết quả. Cấu hình phương pháp gồm checkpoint, kích thước đầu vào, các quy tắc ghép xe với ô và ngưỡng đã lựa chọn trên validation.

Đầu ra của mỗi ảnh gồm mã ô, vị trí vùng ô, trạng thái có xe/trống/chưa xác định và điểm có xe khi phương pháp cung cấp được. Tổng số chỗ trống được tính từ những ô được xác định là trống; số ô chưa xác định được báo riêng. Kết quả được trực quan hóa trên ảnh gốc và lưu dưới dạng CSV/JSON để đánh giá và đối chiếu.

\subsection{Phân biệt phát hiện xe và nhận diện trạng thái ô}
Phát hiện xe trả về vị trí hộp bao của các đối tượng xe và độ tin cậy. Nhận diện trạng thái ô trả lời một vùng ô đỗ đã biết đang có xe hay trống. Với phương pháp dựa trên detector, cần thêm bước ánh xạ hộp xe về ảnh gốc và ghép hình học xe–ô trước khi suy ra trạng thái. Với HOG/SVM, ảnh của từng ô được phân loại trực tiếp theo nhãn trạng thái.

Polygon hoặc hình chữ nhật xoay trong XML PKLot biểu diễn vùng ô đỗ; chúng không xác định hộp bao của chiếc xe. Vì vậy, nhãn ô không được chuyển thành nhãn hộp xe để huấn luyện detector hoặc tính mAP phát hiện xe. Phần đánh giá chính sử dụng trạng thái từng ô, phép đếm và thời gian trên ảnh gốc.

\subsection{Phạm vi và giả định}
Phạm vi chính là bãi xe ngoài trời với camera cố định và các vùng ô đỗ đã được cung cấp. Khi vị trí hoặc góc nhìn camera thay đổi, cấu hình vùng ô phải được cập nhật hoặc xác nhận lại. Hệ thống không giả định một cấu hình polygon duy nhất cho mọi ảnh nếu dữ liệu ghi nhận nhiều biến thể hình học.

Canny và Hough được dùng trên nhánh ảnh xám để phân tích biên, đường kẻ và kiểm tra chất lượng hình học. Nhánh detector sử dụng ảnh màu; việc không tìm thấy đường kẻ không được dùng riêng để kết luận ô trống. Tự động đề xuất vị trí ô đỗ là hướng mở rộng, chưa thuộc phạm vi thực nghiệm chính. Kết quả sẽ được phân tích theo thời tiết và camera; thí nghiệm chuyển góc nhìn trong cùng bãi được trình bày riêng.
